\section{Cumulative transport and the actual Gaussian comparison}
\label{sec:actual-covariance-transport}

We prove the low-frequency transport bounds on the same atlas as
Section~\ref{sec:actual-covariance-fourier-window}. Write $\kappa_i$
for its original real curvature covariance and $\Psi_i$ for the
installed map from level $i$ to level $i+1$, with the identity on
colour coordinates. Installation sums all periodic images of the
cover row. Fix a comparison order $d\ge1$ and a plateau of smoothing
order $J$ with $d-1\le J-2$.

\subsection{Moving differences through an actual step}

For a word $w$ of $j\le d-1$ coordinate directions, the selected
transport has a real quotient row $Q_w$ satisfying
\begin{equation}\label{eq:actual-transport-difference}
 \Delta_w\Psi_i=Q_w\Delta_w,\qquad
 \sup_x\sum_y|Q_w(x,y)|\le R L_i^{-j-2},
\end{equation}
where $L_i=a_i/a_{i+1}$ and $R\ge1$ depends only on the kernel
penalty and adjacent smoothing orders. The quotient needs proof;
a bound on the left side alone would not suffice.

Theorem~\ref{thm:ex-lattice-wire-tower-strip}, applied to the actual
pulse-multiplied curvature column, gives its differentiated row
strip cost $C L_i^{-j-2}$. The original pulse's strip cost simply
multiplies the established column cost. To move coarse differences
to the right, divide the symbol successively by $e^{ip_a}-1$.
The division is regular on the real principal cube before summing
aliases. On the principal alias use
\[
 e^{ip_a/L_i}-1=
 \frac{e^{ip_a}-1}{L_i\operatorname{rail}_{L_i}(p_a/L_i)}.
\]
The denominator has a positive lower bound by the principal rail
estimate. On a nonprincipal alias $q=(p+2\pi l)/L_i$, the column
contains the box filter $A(q)^J$. The geometric-series identity is
\[
 (e^{iq_a}-1)A(q)=(e^{ip_a}-1)s_a(q),
\]
where $s_a$ is a product of the other coordinate box filters and a
phase, divided by $L_i$. For $j$ directions use $j$ box factors.
The remaining power $A^{J-j}$ has nonnegative degree. The remaining
multiplier is continuous by the nonprincipal denominator bounds
already proved. Thus the full differentiated symbol has a continuous
quotient by the product of all coarse differences.

A coefficientwise division also preserves a narrower strip bound.
In one coordinate use the exact finite identity
\[
 \frac{e^{inp_a}-1}{e^{ip_a}-1}
 =\begin{cases}
 \sum_{h=0}^{n-1}e^{ihp_a},&n>0,\\
 -\sum_{h=n}^{-1}e^{ihp_a},&n<0,\\
 0,&n=0.
 \end{cases}
\]
There are $|n|$ terms with no greater displacement. Reducing width
$\sigma$ to $\sigma/2$ costs at most $2/\sigma$, since
$|n|e^{\sigma|n|/2}\le(2/\sigma)e^{\sigma|n|}$.
This divides the symbol minus its value on $p_a=0$; that value
vanishes by the preceding divisibility. Equality off the coordinate
hyperplane extends to it by continuity. Repeat along $w$;
divisibility by the remaining factors survives at each step.
There are finitely many words of length at most $d-1$, so one
narrower strip and one cost work for all. The scalar factor
$L_i^{-j-2}$ is unchanged. Periodization and the triangle inequality
turn this cover bound into an absolute row bound; the sum over the
six two-form types contributes their cardinality. This proves
\eqref{eq:actual-transport-difference} with a cost chosen before $L_i$.

Closedness is preserved on all cochains, including harmonic ones.
Write the actual step as $B_f T H_c$, using reciprocal pulse $B$
and pulse $H$. Its exterior transport $C$ satisfies
\[
 C d_cB_c=d_fB_fT,\qquad B_cH_c=I.
\]
Multiplying by $H_c$ gives $d_f\Psi_i=Cd_c$. No global exactness
of the input cochain is required. Reality follows from the real
cover coefficients.

\subsection{Absolute rows of a product}

For $\ell\le k$ put $m=k-\ell$,
$\mathcal L=\Psi_{k-1}\cdots\Psi_\ell$ and
$S=a_\ell/a_k=\prod_{i=\ell}^{k-1}L_i$.
If every $j$th difference of $f$ is bounded by $v$, iteration gives
\[
       \|\Delta_w\mathcal Lf\|_\infty
                           \le R^mS^{-j-2}v.
\]
This holds for arbitrary real $f$. For a specified output row,
take $f(y)$ to be the sign of $\Delta_w\mathcal L(x,y)$, using
zero when that entry vanishes. Then $|f|\le1$ and every $j$th
difference has absolute value at most $2^j$. Consequently
\begin{equation}\label{eq:actual-cumulative-row}
 S^{j+2}\sum_y|\Delta_w\mathcal L(x,y)|\le R^m2^d.
\end{equation}
The sign test need not be closed. Reality and preservation of
closedness pass separately to the product.

Fix $b\ge6$ and $a>0$. After $R$, choose a plateau block length
$L>1$ so that $4(b+2)\log R\le a\log L$. With
$m=\lceil(b+2)\log H/\log L\rceil$ we have
\[
 R^m\le RH^{a/4}\le H^{a/2},\qquad R^m2^d\le H^a
\]
eventually. The second inequality absorbs fixed $2^d$ into
$H^{a/2}$. Also $S=L^m\ge H^{b+2}$.
The original covariance entries have a fixed bound $C_J$ uniform
in period, by summing the bounded curvature multiplier over its
normalized finite Fourier grid. Paying $C_J\le H^a$, and writing
$P=H^a$, gives
\begin{equation}\label{eq:actual-low-data}
 S^{j+2}\sum_y|\Delta_w\mathcal L(x,y)|\le P,
                  \qquad |\kappa_\ell(y,z)|\le P.
\end{equation}
One finite table of costs is selected before the block length,
then its clearance threshold. On the allowed plateau the prescribed
number of steps is exactly $m$. All estimates therefore concern
the same chosen level and history, uniformly in period.

\subsection{Mean and covariance after normal ordering}

For a real closed coarse cochain $u$ with $\|u\|_\infty\le P$,
set $\mu=\mathcal Lu$ and
$B=\mathcal L\kappa_\ell\mathcal L^{\mathsf T}$.
Finite summation gives
\begin{align}\label{eq:actual-low-moments}
 S^{j+2}|\Delta_w\mu(x)|&\le P^2,\\
 S^{j+2}S^{h+2}|\Delta_w^x\Delta_v^yB(x,y)|&\le P^3.
 \notag
\end{align}
For the second inequality bound each covariance entry by $P$;
the remaining double sum factors into the absolute row sums of
\eqref{eq:actual-low-data}. The mean and both covariance arguments
are closed, by the original closedness and its preservation.

Use $r=H^b$, $a=\lambda/4$, $0<\lambda\le1/4$, and the purchased
condition $S\ge rH^2$. For $j\ge0$,
$r^{j+2}H^4\le S^{j+2}$. Hence the weighted means are at most
$H^{-4}P^2\le H^{-3}$. Retaining only one factor $H^{-4}$ for
the covariance gives the convenient bound
\begin{equation}\label{eq:actual-low-error}
 r^{j+h+4}|\Delta_w^x\Delta_v^yB(x,y)|
                              \le H^{-4+3\lambda/4}.
\end{equation}

Fix distinct continuum positions $x_i$ and their selected rounded
lattice centres $z_i$ at magnification $r$. Required offsets have
subpower size in $H$, so normalized pair displacements stay in
compact sets away from zero. The free covariance and the finitely
many derivatives used here are bounded on these sets.
For an off-diagonal pair, multiply
\eqref{eq:actual-native-model} by $r^{j+h+4}$.
The singular error is at most $CH/r=CH^{1-b}\le H^{-4}/2$
eventually, since $b\ge6$. The period error is bounded by a fixed
power of $H$ divided by the lattice side. The selected mesh estimates
make it at most $H^{-4}/2$ eventually, uniformly in physical period.
Thus the scaled error is at most $H^{-4}$, and the scaled covariance
is bounded by $H^{\lambda/8}$ eventually.

Let $K^0_{ij}$ be the polynomial test-cochain covariance built
from the free jet covariance at $z_i/r$. For both derivative orders
$j,h\le d-2$, at the origins its
weighted differences differ from the corresponding free curvature
derivatives by at most $H^{-4}$. Indeed a difference of order $j$
kills test cochains of degree less than $j$, and the degree-$j$
term recovers its homogeneous curvature tensor exactly. Terms of
degree $q>j$ carry $r^{j-q}\le r^{-1}$. Expanding the finite
covariance matrix in the jet basis bounds the total error by
$C/r\le H^{-4}$ eventually. The coefficient constants are fixed
and the free entries are bounded on the compact configuration set.
If one endpoint has order $d-1$, its derivative of $K^0_{ij}$
is zero, since all test cochains have degree at most $d-2$.
Set $K^0_{ii}=0$, in accordance with free normal ordering.

The actual normal-ordered covariance is
\begin{equation}\label{eq:actual-normalordered-covariance}
 \Gamma_{ij}(s,t)=\mathbf1_{i\ne j}
   \kappa_k(z_i+s,z_j+t)-B(z_i+s,z_j+t).
\end{equation}
Its diagonal is $-B$, generally nonzero. At origins with both
$j,h\le d-2$, the two
comparison errors and \eqref{eq:actual-low-error} give
\[
 r^{j+h+4}|\Delta_w^s\Delta_v^t(\Gamma_{ij}-K^0_{ij})(0,0)|
 \le2H^{-4}+H^{-4+3\lambda/4}\le H^{-3}
\]
eventually; only the last term is needed on the diagonal. For $d=1$
this range of origin derivatives is empty.
At a top-order endpoint the polynomial comparison vanishes and
\[
 H^{\lambda/8}+H^{-4+3\lambda/4}\le H^{\lambda/4}=P.
\]
These, with the mean and closedness bounds, are precisely the
finite-reading estimates with $P=H^{\lambda/4}$ and $E=H^{-3}$
used in Section~\ref{sec:allarity-coordinate-calibration}.
They hold uniformly over the admitted closed real coarse backgrounds.
Positive semidefiniteness of $\Gamma$ is not assumed.

Application must retain the purchased plateau and clearance conditions
and the identification of this atlas with the actual reader law.
These belong to the common history and reader construction; an
arbitrary symmetric model cannot replace that identification.
